% ============================================================================
% Fichier  : partie4/P04_C16_processus_raci_mesure.tex
% Titre    : Processus : Modéliser, Responsabiliser, Mesurer
% Auteur   : MINKA MI NGUIDJOÏ Thierry Emmanuel
% Version  : 1.0 -- octobre 2026
% ============================================================================
% Sources : BPMN 2.0 (OMG, ISO/IEC 19510) ; COBIT 2019 (dossier Normes) ;
% ITIL 4 (sources publiques). La correspondance avec les articles de
% l'ISO/IEC 20000-1:2018 est établie par nos soins (norme absente du
% dossier Normes) : A VERIFIER dans le texte de la norme.
% Laboratoire associé : GCI-LAB-09 (matrice RACI et correspondance COBIT-ITIL
% d'un processus de NOUN PAY).
% ============================================================================

\chapter{Processus : Modéliser, Responsabiliser, Mesurer}
\label{chap:processus}
\chaptermark{Processus : modéliser, responsabiliser}

\epigraph{%
    « Quand tout le monde est responsable, personne ne l'est. »%
}{Adage du management}

\niveauChapitre{Expert}{%
    Maîtrise des Chapitres~\ref{chap:cobit-objectifs} (les objectifs COBIT)
    et~\ref{chap:itil-niveaux} (niveaux de service et fournisseurs).%
}

\begin{boxobjectifs}
\begin{itemize}
    \item Modéliser un processus en BPMN.
    \item Construire une matrice RACI sans ambiguïté.
    \item Mettre en correspondance COBIT, ITIL et ISO/IEC~20000-1.
    \item Définir les indicateurs (KPI, KGI) et un tableau de bord.
    \item Désigner les propriétaires de processus et organiser
          l'amélioration continue.
\end{itemize}
\end{boxobjectifs}

\bigskip

% ============================================================================
\section*{Question directrice : qui fait quoi, et comment le sait-on~?}
% ============================================================================

Les Chapitres~\ref{chap:itil-svs} à~\ref{chap:itil-niveaux} ont présenté
les pratiques une par une. Mais une organisation ne vit pas des pratiques
isolées~: elle vit des enchaînements d'activités, qui traversent les
équipes, les fournisseurs et les outils. Chez NOUN PAY, à Bafoussam, la
mise à jour du modèle de notation de crédit passe par cinq personnes et un
fournisseur de nuage~; personne ne sait qui l'autorise. Chez KRIBILOG,
NETPORT SERVICES et le centre de services se renvoient les incidents.

Ce chapitre apprend à assembler les pratiques en processus maîtrisés~: les
modéliser (section~\ref{sec:c16-bpmn}), répartir les responsabilités
(section~\ref{sec:c16-raci}), établir la correspondance entre les
référentiels (section~\ref{sec:c16-correspondance}), mesurer
(section~\ref{sec:c16-mesure}) et améliorer (section~\ref{sec:c16-proprietaires}).
Le laboratoire GCI-LAB-09 fait construire la matrice RACI et la
correspondance d'un processus de NOUN PAY.

\separateur

% ============================================================================
\section{Processus et Pratiques}
\label{sec:c16-notions}
% ============================================================================

\begin{boxdef}
\termecle{Processus}\index{processus}~: ensemble d'activités corrélées qui
transforme des entrées en sorties, en vue d'un objectif. Il a un début, une
fin, des acteurs et des résultats mesurables.

\medskip
\termecle{Propriétaire de processus}\index{propriétaire de processus}
(\textit{process owner})~: personne qui répond de la conception, de la
performance et de l'amélioration d'un processus dans toute
l'organisation. Elle ne l'exécute pas nécessairement elle-même.

\medskip
\termecle{Gestionnaire de processus}\index{gestionnaire de processus}
(\textit{process manager})~: personne qui pilote l'exécution quotidienne du
processus, sur un périmètre donné.
\end{boxdef}

Un processus n'est pas une pratique. Une \emph{pratique} ITIL ou un
\emph{objectif} COBIT décrivent un domaine de maîtrise, avec toutes ses
composantes (Chapitre~\ref{chap:cobit-archi}). Un \emph{processus} est
l'enchaînement concret par lequel une organisation réalise une partie de
ce domaine. La pratique « gestion des incidents » peut donner lieu à
plusieurs processus~: incidents ordinaires, incidents majeurs, incidents
de sécurité.

% ============================================================================
\section{Modéliser en BPMN}
\label{sec:c16-bpmn}
% ============================================================================

\subsection{Les éléments du langage}

Le langage BPMN (\textit{Business Process Model and Notation}), publié par
l'Object Management Group dans sa version~2.0 et repris par la norme
ISO/IEC~19510, est la notation de référence pour modéliser les
processus. Il suffit de quelques éléments pour en tirer l'essentiel.

\begin{tableaucours}{L{3.0cm}L{3.6cm}L{5.6cm}}{%
    \tableauHeader{Élément} &
    \tableauHeader{Représentation} &
    \tableauHeader{Usage}
}
Événement & Cercle (début, intermédiaire, fin) & Ce qui déclenche, ponctue
ou termine le processus \\
Activité & Rectangle arrondi & Une tâche ou un sous-processus \\
Passerelle & Losange & Un choix (exclusif), une séparation en parallèle,
une jonction \\
Flux de séquence & Flèche pleine & L'ordre des activités \\
Flux de messages & Flèche en pointillé & Un échange entre deux acteurs \\
Couloir & Bande horizontale & Le rôle ou l'organisation qui exécute les
activités \\
\end{tableaucours}
\vspace{-0.8em}
\begin{center}{\small\itshape Tableau~16.1. Les éléments essentiels de
BPMN~2.0}\end{center}

\subsection{Un exemple}

\noindent\begin{minipage}{\linewidth}
\begin{center}
\begin{tikzpicture}[
    tache/.style = {
        draw=CouleurBordDef, fill=CouleurDef, rounded corners=3pt,
        text width=1.8cm, minimum height=0.9cm, align=center,
        font=\scriptsize\color{CouleurPrimaire}
    },
    evt/.style = {circle, draw=CouleurPrimaire, thick, minimum size=0.5cm},
    evtfin/.style = {circle, draw=CouleurPrimaire, line width=2pt,
        minimum size=0.5cm},
    pass/.style = {diamond, draw=CouleurAccent, fill=CouleurAccent!15,
        minimum size=0.8cm, font=\scriptsize\bfseries\color{CouleurPrimaire},
        inner sep=0pt},
    fleche/.style = {-{Stealth[length=4pt]}, thick, color=CouleurSecondaire},
    couloir/.style = {font=\scriptsize\bfseries\color{CouleurPrimaire},
        rotate=90, anchor=south}
]
% Couloirs
\draw[CouleurPrimaire] (-0.6,0.9) rectangle (13.0,-1.1);
\draw[CouleurPrimaire] (-0.6,-1.1) rectangle (13.0,-3.1);
\node[anchor=north west, font=\tiny\bfseries\color{CouleurPrimaire}]
    at (-0.6,0.9) {Centre de services (KRIBILOG)};
\node[anchor=north west, font=\tiny\bfseries\color{CouleurPrimaire}]
    at (-0.6,-1.1) {NETPORT SERVICES (infogérant)};
% Centre de services
\node[evt] (d) at (0.2,-0.1) {};
\node[tache] (e) at (1.8,-0.1) {Enregistrer et prioriser};
\node[pass] (p) at (3.7,-0.1) {X};
\node[tache] (r1) at (5.5,-0.1) {Résoudre au niveau~1};
\node[tache] (v) at (10.2,-0.1) {Vérifier avec l'utilisateur};
\node[evtfin] (f) at (12.2,-0.1) {};
% NETPORT
\node[tache] (dg) at (5.5,-2.1) {Diagnostiquer};
\node[tache] (r2) at (8.0,-2.1) {Résoudre ou contourner};
% Flux
\draw[fleche] (d) -- (e);
\draw[fleche] (e) -- (p);
\draw[fleche] (p) -- node[above, font=\tiny]{connu} (r1);
\draw[fleche] (p) |- node[pos=0.2, right, font=\tiny]{inconnu} (dg);
\draw[fleche] (dg) -- (r2);
\draw[fleche] (r2) -| (v);
\draw[fleche] (r1) -- (v);
\draw[fleche] (v) -- (f);
\end{tikzpicture}
\end{center}
\vspace{-0.4em}
\begin{center}{\small\itshape Figure~16.1. Processus de traitement d'un
incident chez KRIBILOG (BPMN simplifié)}\end{center}
\end{minipage}

\medskip

La figure fait apparaître ce que le texte des procédures masquait~: le
passage d'un couloir à l'autre. C'est à ces frontières que se perdent les
incidents. Chez KRIBILOG, la flèche de retour de NETPORT SERVICES vers la
vérification n'existait pas~: l'infogérant clôturait lui-même les tickets
(Chapitre~\ref{chap:itil-incidents}).

\begin{boxalert}
\textbf{Violation typique.} Une organisation modélise ses processus tels
qu'ils \emph{devraient} être, à partir des référentiels, sans observer ce
qui se passe réellement. Les schémas sont beaux, mais personne ne les
applique.

\smallskip
\textbf{Règle~:} modéliser d'abord le processus \emph{tel qu'il est}, à
partir d'entretiens et des traces de l'outil (principe « commencer là où
l'on est », Chapitre~\ref{chap:itil-svs}), puis le processus cible, et
mesurer l'écart.
\end{boxalert}

% ============================================================================
\section{Responsabiliser : la Matrice RACI}
\label{sec:c16-raci}
% ============================================================================

\subsection{Quatre rôles}

\begin{boxdef}
\termecle{Matrice RACI}\index{matrice RACI}~: tableau qui croise les
activités d'un processus et les rôles de l'organisation, et indique pour
chaque case le rôle tenu~:
\begin{itemize}
    \item \textbf{R} (\textit{Responsible})~: réalise l'activité~;
    \item \textbf{A} (\textit{Accountable})~: en répond et l'approuve~;
    \item \textbf{C} (\textit{Consulted})~: est consulté avant~;
    \item \textbf{I} (\textit{Informed})~: est informé après.
\end{itemize}
\end{boxdef}

COBIT~2019 fournit pour chaque pratique une matrice RACI type, établie
avec des rôles génériques (Chapitre~\ref{chap:cobit-objectifs})~; il
revient à chaque organisation de la traduire avec ses propres fonctions.

\subsection{Cinq règles}

\begin{enumerate}
    \item \textbf{Un seul A par activité.} Deux personnes qui « répondent »
          d'une même activité, c'est n'en avoir aucune.
    \item \textbf{Au moins un R par activité.} Une activité sans exécutant
          ne sera pas faite.
    \item \textbf{Le A au bon niveau.} Celui qui répond doit avoir
          l'autorité de décider et les moyens d'agir.
    \item \textbf{Peu de C.} Trop de consultations ralentissent tout~;
          consulter, c'est s'engager à attendre la réponse.
    \item \textbf{Des noms de fonctions, pas de personnes.} La matrice
          survit aux départs.
\end{enumerate}

\begin{boxscenario}{Opération NOUN. Qui autorise~?}
\textbf{Faits.} Chez NOUN PAY, une nouvelle version du modèle de notation
de crédit est préparée par l'équipe de science des données, testée par le
directeur technique, déployée par l'équipe DevOps sur l'infrastructure du
fournisseur de nuage. La fondatrice est « tenue au courant ». Le délégué à
la protection des données n'est pas associé.

\medskip
\textbf{Matrice constatée.}

\medskip
\begin{tableaucours}{L{3.4cm}ccccc}{%
    \tableauHeader{Activité} &
    \tableauHeader{DG} &
    \tableauHeader{Dir. tech.} &
    \tableauHeader{Données} &
    \tableauHeader{DevOps} &
    \tableauHeader{DPD}
}
Préparer la version & I & A & R & · & · \\
Tester & · & A/R & C & · & · \\
Évaluer l'impact sur les clients & · & · & · & · & · \\
Autoriser la mise en production & I & A & · & A & · \\
Déployer & · & · & · & A/R & · \\
\end{tableaucours}

\medskip
\textbf{Constatation.} L'activité « évaluer l'impact sur les clients » n'a
ni R ni A. L'autorisation compte deux A. Le délégué à la protection des
données n'apparaît nulle part.

\medskip
\textbf{Analyse} (confiance élevée). Trois règles sur cinq sont violées.
C'est l'explication directe de la demande de l'Autorité de protection des
données (Chapitre~\ref{chap:droit})~: personne n'était chargé de vérifier
qu'une nouvelle version ne pénalisait pas une catégorie de clients.
\end{boxscenario}

\begin{boxprocedure}
\textbf{Une matrice RACI auditable.} Pour qu'une matrice RACI soit
opposable, elle doit~:
\begin{itemize}
    \item porter sur des activités précises, rattachées à une pratique ou à
          un objectif identifié~;
    \item utiliser les fonctions réelles de l'organisation, telles qu'elles
          figurent à l'organigramme~;
    \item respecter les règles « un seul A » et « au moins un R »~;
    \item être approuvée par le propriétaire du processus et datée~;
    \item être cohérente avec les fiches de poste et les délégations de
          pouvoir.
\end{itemize}
\end{boxprocedure}

% ============================================================================
\section{COBIT, ITIL et ISO/IEC 20000-1}
\label{sec:c16-correspondance}
% ============================================================================

\subsection{Pourquoi une correspondance}

Une organisation qui utilise COBIT pour sa gouvernance, ITIL pour ses
pratiques et vise l'ISO/IEC~20000-1 doit éviter de construire trois
dispositifs parallèles. La correspondance permet de n'avoir qu'un
processus, documenté une fois, qui répond aux trois référentiels.

\begin{tableaucours}{L{3.0cm}L{4.6cm}L{4.6cm}}{%
    \tableauHeader{COBIT 2019} &
    \tableauHeader{Pratiques ITIL 4} &
    \tableauHeader{ISO/IEC 20000-1:2018}
}
APO08 Relations & Gestion des relations & 8.3.2 Relations avec les
clients \\
APO09 Accords de service & Gestion des niveaux de service, du catalogue &
8.3.3 Niveaux de service~; 8.2.4 Catalogue \\
APO10 Fournisseurs & Gestion des fournisseurs & 8.3.4 Fournisseurs \\
APO06 Budget et coûts & Gestion financière des services & 8.4.1
Budgétisation et comptabilité \\
BAI04 Disponibilité et capacité & Gestion de la disponibilité, de la
capacité et des performances & 8.7.1 Disponibilité~; 8.4.3 Capacité \\
BAI06 Changements & Habilitation des changements & 8.5.1 Changements \\
BAI07 Transition & Mises en production, déploiements, validation et tests
& 8.5.2 Conception et transition~; 8.5.3 Mises en production et
déploiement \\
BAI09 Actifs & Gestion des actifs informatiques & 8.2.5 Actifs \\
BAI10 Configuration & Gestion de la configuration des services & 8.2.6
Configuration \\
DSS02 Demandes et incidents & Centre de services, incidents, demandes de
service & 8.6.1 Incidents~; 8.6.2 Demandes de service \\
DSS03 Problèmes & Gestion des problèmes & 8.6.3 Problèmes \\
DSS04 Continuité & Gestion de la continuité des services & 8.7.2
Continuité \\
APO13, DSS05 Sécurité & Gestion de la sécurité de l'information & 8.7.3
Sécurité de l'information \\
MEA01 Performance & Mesure et compte rendu & 9.1 Surveillance et mesure \\
\end{tableaucours}
\vspace{-0.8em}
\begin{center}{\small\itshape Tableau~16.2. Correspondance entre COBIT~2019,
ITIL~4 et l'ISO/IEC~20000-1:2018 (établie par nos soins~; numéros
d'articles de la norme à vérifier dans son texte)}\end{center}

\subsection{Les limites d'une correspondance}

Une correspondance n'est jamais une équivalence. Trois précautions
s'imposent.
\begin{itemize}
    \item \textbf{Les périmètres diffèrent.} L'objectif BAI07 de COBIT couvre
          la conversion des données et la revue après mise en œuvre, que la
          pratique ITIL de gestion des mises en production ne traite pas de
          la même manière.
    \item \textbf{Les natures diffèrent.} COBIT fixe un objectif, ITIL décrit
          une pratique, la norme formule une exigence~: satisfaire l'une ne
          garantit pas de satisfaire les autres.
    \item \textbf{Les éditions changent.} Une correspondance se date et se
          révise à chaque nouvelle édition d'un des référentiels.
\end{itemize}

\begin{boxmetier}{Responsable GRC}
\textbf{Construire un référentiel unique de processus}
\begin{enumerate}
    \item Partir des objectifs COBIT prioritaires
          (Chapitre~\ref{chap:cobit-conception}).
    \item Pour chacun, désigner le processus de l'organisation qui le
          réalise, et la pratique ITIL qui guide sa conception.
    \item Repérer les exigences de l'ISO/IEC~20000-1 si une certification
          est visée, et les obligations légales (Chapitre~\ref{chap:droit}).
    \item Documenter le processus une seule fois, avec un tableau de
          correspondance en annexe.
    \item Faire approuver l'ensemble par le comité stratégique des SI.
\end{enumerate}
\end{boxmetier}

% ============================================================================
\section{Mesurer : KPI, KGI et Tableau de Bord}
\label{sec:c16-mesure}
% ============================================================================

\subsection{Deux familles d'indicateurs}

\begin{boxdef}
\termecle{Indicateur clé d'objectif}\index{indicateur clé d'objectif}
(\textit{key goal indicator}, \sigle{KGI})~: indicateur qui dit si un
objectif est atteint. Il mesure un résultat, après coup.

\medskip
\termecle{Indicateur clé de performance}\index{indicateur clé de
performance} (\textit{key performance indicator}, \sigle{KPI})~:
indicateur qui dit si un processus fonctionne bien et a des chances
d'atteindre l'objectif. Il mesure une performance, en cours de route.
\end{boxdef}

Cette distinction, mise en avant par COBIT~4.1~\cite{itgi2007cobit41}, reste
utile~: COBIT~2019 propose pour chaque objectif d'entreprise, chaque
objectif d'alignement et chaque pratique des « exemples
d'indicateurs », qui relèvent de l'une ou l'autre famille.

\begin{tableaucours}{L{3.0cm}L{4.6cm}L{4.6cm}}{%
    \tableauHeader{Niveau} &
    \tableauHeader{Indicateur (KRIBILOG)} &
    \tableauHeader{Destinataire}
}
Objectif d'entreprise (KGI) & Nombre de navires retardés par une panne du
suivi & Conseil d'administration \\
Objectif d'alignement (KGI) & Disponibilité de bout en bout du service de
suivi & Comité stratégique des SI \\
Processus (KPI) & Taux de réouverture des incidents~; âge moyen des
problèmes ouverts & Responsable des services informatiques \\
Activité (KPI) & Part des tickets enregistrés dans les 5~minutes &
Responsable du centre de services \\
\end{tableaucours}
\vspace{-0.8em}
\begin{center}{\small\itshape Tableau~16.3. La cascade des indicateurs}\end{center}

\subsection{Un bon indicateur}

Un indicateur utile est \textbf{défini} (formule, source, fréquence),
\textbf{assorti d'une cible}, \textbf{attribué} à un responsable,
\textbf{difficile à manipuler} (accompagné de son indicateur de
contrôle, Chapitre~\ref{chap:itil-incidents}) et \textbf{suivi d'action}~:
un indicateur dont l'écart ne déclenche jamais rien est un indicateur
inutile.

\begin{boxobjet}{MEA01 · Gérer la surveillance de la performance et de la conformité}
\textbf{Finalité~:} assurer la transparence de la performance et de la
conformité, et piloter l'atteinte des objectifs.

\textbf{Contenu (résumé)~:} fixer l'approche de surveillance, définir les
cibles, collecter et traiter les données, analyser et rendre compte de la
performance, assurer la mise en œuvre des actions correctives.

\textbf{Sorties~:} tableau de bord, rapports de performance, plans d'action.

\textbf{RACI (extrait, KRIBILOG)~:} \emph{A} au directeur général,
\emph{R} au responsable des services informatiques et au responsable GRC,
\emph{I} au conseil.

\textbf{Indicateurs~:} part des indicateurs dotés d'une cible et d'un
responsable~; délai entre l'écart constaté et l'action décidée.
\end{boxobjet}
{\small\itshape Résumé établi par nos soins à partir de
COBIT~2019~\cite{isaca2018cobitobjectifs}.}

% ============================================================================
\section{Propriétaires et Amélioration Continue}
\label{sec:c16-proprietaires}
% ============================================================================

Un processus sans propriétaire se dégrade. Le propriétaire répond de
quatre choses~: la conception du processus et sa documentation, la
définition de ses indicateurs, sa performance dans toute l'organisation,
et son amélioration. Il conduit, au moins une fois par an, une revue du
processus~: indicateurs, incidents, constats d'audit, retours des
utilisateurs, et alimente le registre des améliorations
(Chapitre~\ref{chap:itil-svs}).

\begin{boxmetier}{DSI}
\textbf{Désigner les propriétaires de processus}
\begin{enumerate}
    \item Un propriétaire par processus prioritaire, de rang suffisant pour
          imposer ses décisions aux équipes concernées.
    \item Pour les processus qui traversent un fournisseur, un propriétaire
          \emph{interne}~: on n'externalise pas la propriété d'un processus.
    \item Une lettre de mission, qui précise les pouvoirs, les indicateurs et
          le rythme de la revue.
    \item Une présentation annuelle de chaque propriétaire au comité
          stratégique des SI.
\end{enumerate}
\end{boxmetier}

\begin{boiteRemarque}{Le processus n'est-il pas l'ennemi de l'agilité~?}
Un processus lourd, imposé d'en haut et jamais révisé, l'est sans doute.
Mais l'absence de processus n'est pas l'agilité~: c'est la dépendance aux
personnes. Les organisations agiles ont des processus courts, visibles,
automatisés quand c'est possible, et révisés souvent~; elles appliquent les
principes d'ITIL « rester simple et pratique » et « progresser par
itérations ». Le Chapitre~\ref{chap:agile-cloud} y reviendra.
\end{boiteRemarque}

% ============================================================================
\separateur
% ============================================================================

\begin{boxresume}
\textbf{Ce qu'il faut retenir de ce chapitre~:}
\begin{itemize}
    \item \textbf{Processus et pratique}~: la pratique décrit un domaine de
          maîtrise, le processus l'enchaînement concret par lequel
          l'organisation le réalise.
    \item \textbf{BPMN} rend visibles les frontières entre acteurs, là où se
          perdent les dossiers. Modéliser d'abord l'existant.
    \item \textbf{RACI}~: un seul A, au moins un R, le A au bon niveau, peu
          de C, des fonctions et non des personnes.
    \item \textbf{La correspondance COBIT, ITIL, ISO/IEC~20000-1} permet un
          seul processus pour trois référentiels~; ce n'est jamais une
          équivalence.
    \item \textbf{KGI et KPI}~: le premier dit si l'objectif est atteint, le
          second si le processus a des chances de l'atteindre. Les
          indicateurs se déclinent du conseil à l'activité.
    \item \textbf{Un propriétaire par processus}, interne même quand le
          processus traverse un fournisseur, qui en conduit la revue
          annuelle.
\end{itemize}
\end{boxresume}

\bigskip

\begin{boxexo}[title={\ding{45}\quad Exercice~16.1. Corriger une matrice RACI \niveaubase}]
Reprenez la matrice constatée de NOUN PAY (section~\ref{sec:c16-raci}).
\begin{enumerate}[(a)]
    \item Listez toutes les violations des cinq règles.
    \item Proposez une matrice corrigée, en ajoutant les activités qui
          manquent.
    \item Indiquez qui doit approuver la matrice corrigée.
\end{enumerate}
\end{boxexo}

\begin{boxexo}[title={\ding{45}\quad Exercice~16.2. Modéliser le processus de changement de l'ANPP \niveaumoyen}]
Modélisez en BPMN le processus de changement normal de l'ANPP, depuis la
demande d'une direction métier jusqu'à la clôture, avec quatre couloirs~:
direction métier, division informatique, comité des changements,
prestataire.
\begin{enumerate}[(a)]
    \item Faites figurer au moins une passerelle exclusive (autorisé ou
          refusé) et un flux de messages.
    \item Repérez les frontières où le dossier risque de se perdre.
    \item Établissez la matrice RACI correspondante.
    \item Rattachez chaque activité à une pratique de BAI06 et à l'article
          correspondant de l'ISO/IEC~20000-1.
\end{enumerate}
\end{boxexo}

\begin{boxexo}[title={\ding{45}\quad Exercice~16.3. Le tableau de bord de la BANQUE SANAGA \niveauexpert}]
Construisez le tableau de bord de la gouvernance des SI de la BANQUE
SANAGA pour l'année qui suit la panne.
\begin{enumerate}[(a)]
    \item Proposez deux KGI par niveau (objectifs d'entreprise prioritaires,
          objectifs d'alignement) et deux KPI pour chacun des dix objectifs
          prioritaires (Chapitre~\ref{chap:cobit-objectifs}).
    \item Pour chaque indicateur, précisez la formule, la source, la
          fréquence, la cible et le responsable.
    \item Associez à chaque indicateur susceptible d'être manipulé son
          indicateur de contrôle.
    \item Indiquez ce qui sera présenté au conseil, au comité stratégique
          des SI et au comité d'audit.
\end{enumerate}
\end{boxexo}

\bigskip

\begin{boxinfo}
\textbf{Transition.} Le cours a désormais présenté toutes les pièces~: la
gouvernance, COBIT, le cadre normatif et juridique, ITIL et les
processus. Le Chapitre~\ref{chap:cas-integrateurs} les assemble sur les
quatre cas fil rouge, en études de cas intégratrices.
\end{boxinfo}
